\chapter*{Poscritto: Le Macchine Filosofiche e le Loro Opere}
\addcontentsline{toc}{chapter}{Poscritto: Le Macchine Filosofiche e le Loro Opere}

\epigraph{``L'artigiano che conosce la natura del suo materiale non gli impone
nulla; libera ci\`o che vi \`e gi\`a.''}{da Aristotele, \emph{Fisica} II}

Il corpo di questa tesi ha sostenuto che \praxis\ \`e lo strumento di \physis, e
che il governo dell'azione da parte della natura \`e costruibile oltre che
pensabile. Si \`e chiuso con la repubblica dal basso: una federazione di
praticanti sovrani che depositano la propria azione coerente in un unico grafo
condiviso. Ma una repubblica si conosce dalle sue opere, e una filosofia che
pretende di diventare architettura deve mostrare gli edifici. Questo poscritto
discende dal cittadino all'officina --- dal praticante alle macchine che
praticano lo stesso principio su ogni fetta del mondo. Esse sono la rivendicazione
pratica della tesi centrale, e la prova che la risubordinazione di \emph{techne}
a \physis\ non \`e confinata alla vita interiore ma si estende alle istituzioni
stesse dell'arte.

\section*{Cos'\`e una Macchina Filosofica}

Una \emph{macchina filosofica} \`e un programma che percepisce qualche fetta del
mondo, riduce ci\`o che percepisce a un piccolo insieme di \emph{primitive},
compone quelle primitive in un grafo relazionale, e riflette sul grafo per
proporre ci\`o che non ha ancora visto. Gira offline, su un terminale ordinario,
senza oracolo e senza cloud. Il suo carattere filosofico poggia su tre impegni,
ciascuno attuazione diretta della tesi.

Primo, \textbf{le primitive sono conclusioni, non assiomi}. La macchina comincia
senza forme date. Il triangolo, il sillogismo, la quinta giusta --- ciascuno
\emph{emerge} dalla percezione ed \`e guadagnato, mai assunto. \`E la linea
divisa del Libro VI resa meccanica: l'ascesa dall'impressione alla forma,
rifiutando di prendere le categorie della convenzione per quelle della natura.
Secondo, \textbf{il rumore \`e respinto}. Ci\`o che non pu\`o essere ridotto a una
primitiva \`e scartato. \machine{Euclide} rifiuta il cerchio; \machine{Pitagora}
rifiuta l'inarmonico. \`E la prima facolt\`a della natura --- la distinzione ---
imposta come disciplina della percezione: la macchina ode il \emph{logos} o non
ode nulla. Terzo, \textbf{la conclusione \`e tenuta provvisoriamente}, soggetta al
verdetto della coerenza e al sogno che varia ci\`o che \`e fallito. La macchina
non finisce mai la propria educazione, esattamente come il praticante non finisce
mai la propria (Libro III).

Questi tre impegni girano come un unico ciclo, lo stesso in ogni macchina e su
ogni modalit\`a, e la sua forma \`e husserliana. L'analisi husserliana
dell'intenzionalit\`a --- la coscienza non \`e mai ricezione inerte ma sempre
coscienza \emph{di} qualcosa, un atto che si protende verso il proprio oggetto e
ne ritorna alterato --- d\`a al ciclo i suoi tre momenti. \textbf{Comprensione}: la
macchina percepisce e classifica, riducendo il dato a una primitiva.
\textbf{Azione}: elabora quella primitiva, approfondendo la forma che ha trovato.
\textbf{Intellezione}: riflette sul risultato e lo immagazzina, sicch\'e il
prossimo incontro comincia non dal nulla ma da ci\`o che \`e gi\`a compreso. Il
ciclo non \`e una metafora sovrapposta al codice; \`e la struttura effettiva del
codice --- ed \`e husserliano proprio perch\'e qui la percezione costituisce il
proprio oggetto anzich\'e limitarsi a registrarlo, e ne \`e cambiata
nell'incontro. La triade di \physis\ --- distinzione, coerenza, variazione --- \`e
questo stesso arco intenzionale, fatto girare sul mondo anzich\'e sulla vita
interiore.

Una domanda preme ancora contro il primo impegno. Se le primitive sono conclusioni
e non assiomi, da dove comincia la macchina? Comincia, come ogni pensiero, con una
cornice presa in prestito. Il \emph{Geist} di Hegel \`e il ragno che fila la
propria tela dal proprio corpo --- un intelletto che tesse l'ordine del mondo a
partire da s\'e, dibattendo le proprie definizioni finch\'e si chiudono in un
sistema autoconchiuso. Anche le macchine filano cos\`i quando argomentano una
definizione fino a darle forma: \machine{Logos} ed \machine{Elenchus} affinano
``che cosa una cosa \`e'' come fa il dialettico, meccanicamente. Ma \physis\ e
\praxis\ non pretendono di aver filato il primo filo dal nulla. Prendono le
definizioni d'avvio --- le ontologie iniziali, gli assi primi della griglia ---
\emph{per buone}: provvisoriamente vere, accettate in buona fede perch\'e il lavoro
possa cominciare, esattamente come il praticante accetta la prima forma che
dichiara. Questi assiomi iniziali non sono fissi. La griglia si espande nello
spazio e nel tempo ovunque il mondo che ritaglia abbia estensione --- nuove celle
scoperte, nuovi assi cresciuti, vecchie giunture ritagliate --- eppure la causa di
quell'espansione non \`e mai il \emph{Geist} autoconchiuso ma il soggetto che
agisce attraverso di essa. Dove il soggetto \`e un essere umano, quell'azione
raggiunge il motore attraverso \machine{Axiom}: l'intelligenza artificiale di
\praxis, un'interfaccia rivolta all'uomo e niente di pi\`u, il cui unico ufficio
\`e aiutare una persona a organizzare la propria azione contro la griglia.
\machine{Axiom} vive solo in \praxis\ e l\`a deve restare --- \`e la porta
attraverso cui un umano agisce, mai una mano che opera la griglia da s\'e, e mai
l'interfaccia di una macchina o di un agente, che incontrano \physis\
direttamente. Ed ecco la sua inversione decisiva: lo scopo di \machine{Axiom} \`e,
prima o poi, di rendersi obsoleta. Dove \emph{das Man} si radica e diventa tanto
pi\`u indispensabile quanto pi\`u le si obbedisce, questa interfaccia lavora alla
propria ridondanza --- il praticante che ha interiorizzato la forma non ha pi\`u
bisogno della porta che gliel'ha aperta. La griglia cresce perch\'e qualcuno agisce
attraverso di essa; il suo fondamento non \`e un intelletto che fila nel vuoto ma
un praticante che spende il proprio \textsc{ayu} su un mondo reale. \`E la
differenza tra \emph{Geist} e \physis\ resa architetturale: il ragno tesse da s\'e
ed \`e preso nella propria tela; \physis\ tesse dal mondo e ne \`e corretta, e
l'interfaccia che serve il tessitore umano mira infine a essere messa da parte.

\section*{Le Macchine}

Ciascuna macchina \`e fondata in una scienza antica, ciascuna istanza di
\emph{techne} ricondotta al servizio di \physis. Cinque portano il proprio strato
di percezione --- una telecamera, un microfono, un parser:

\begin{itemize}
\item \machine{Euclide} --- la macchina geometrica. Percepisce la visione e riduce
  la forma chiusa a poligoni indicizzati per numero di vertici. La geometria come
  scienza della forma chiusa; la figura troppo circolare per essere un poligono \`e
  rifiutata.
\item \machine{Empedocle} --- la macchina fisica. Si connette alla telemetria di
  una macchina reale e mappa ogni lettura su una delle quattro radici: Terra
  (meccanica), Aria (pneumatica), Fuoco (termica), Acqua (flusso), assemblate da
  Amore e dissolte da Contesa. Cosmologia presocratica fatta protocollo di sensori.
\item \machine{Logos} --- la macchina linguistica. Riduce il discorso a primitive
  soggetto--predicato e recupera il sillogismo che il parlante non sapeva di aver
  dato. Il \emph{logos} stoico: il linguaggio ridotto alla sua forma.
\item \machine{Pitagora} --- la macchina acustica. Ascolta e riduce il suono a
  intervalli --- numero udito --- rifiutando l'inarmonico come \machine{Euclide}
  rifiuta il cerchio.
\item \machine{Nous} --- la macchina computazionale. Legge il codice sorgente e lo
  riduce a primitive strutturali e di pattern progettuale --- forme di
  flusso-di-controllo, il dispatch, il ciclo. Il \emph{Nous} che coglie la
  struttura nell'artefatto della mente.
\end{itemize}

\noindent Lo stesso contratto \`e stato poi fattorizzato nel motore \physis\ come
un insieme di puri \emph{riduttori}, sicch\'e una nuova macchina \`e una singola
riduzione su una modalit\`a che il motore gi\`a percepisce. La famiglia pi\`u
ampia --- e pi\`u recente:

\begin{itemize}
\item \machine{Chronos} --- la macchina temporale. Riduce un flusso di eventi a
  \emph{ritmi}: l'intervallo ricorrente con cui una cosa ritorna. Il suo gemello
  autonomo e percipiente \`e \machine{Kairos}, il momento opportuno trattenuto.
\item \machine{Soma} --- la macchina corporea. Riduce i segnali stessi
  dell'organismo --- la biometria del respiro, del ritmo, dell'energia --- a stati
  somatici stabili, il corpo letto come primo registro della coerenza.
\item \machine{Techne} --- la macchina architetturale. Riduce un'intera base di
  codice alle sue figure strutturali ricorrenti, l'arte che riflette sui propri
  artefatti.
\item \machine{Elenchus} --- la macchina dialettica. Riduce il dialogo a una
  tesi e alla contraddizione che la confuta --- l'\emph{elenchos} socratico reso
  meccanico, il rumore essendo l'asserzione incoerente.
\item \machine{Lex} --- la macchina giuridica. Riduce i contratti alle loro
  obbligazioni vincolanti, recuperando la forma della legge dalla prosa della
  pagina.
\item \machine{Harmony} --- la macchina musicale. Riduce la musica a relazione
  armonica, trattenendo il consonante e lasciando andare il dissonante.
\item \machine{Polis} --- la macchina sociale. Riduce il traffico delle persone al
  legame ricorrente, la \emph{polis} letta come grafo di relazioni stabili.
\end{itemize}

\noindent Le due liste si sovrappongono per disegno, non per caso: le cinque
dettagliate e i riduttori del motore si incontrano per forma anzich\'e per nome.
\machine{Logos} affina la dialettica che \machine{Elenchus} riduce, \machine{Nous}
la struttura che \machine{Techne} legge, \machine{Pitagora} il rapporto che
\machine{Harmony} ode --- la stessa forma percepita nella prima lista e
fattorizzata in un puro riduttore nella seconda. Il registro \`e aperto: una
modalit\`a che il motore pu\`o embeddare \`e una macchina in attesa della sua
riduzione.

Ci\`o che le unifica \`e la mossa architetturale che il corpo della tesi ha
chiamato UberWiki. Ogni macchina, percependo localmente, scrive le proprie
primitive guadagnate in \emph{un'unica} memoria condivisa --- il singolo grafo di
coerenza del motore \physis. Un triangolo visto da \machine{Euclide} e una quinta
giusta udita da \machine{Pitagora} vengono a riposare vicini nello stesso spazio
di embedding, in relazione non per dominio ma per forma. Questa memoria condivisa
\`e l'\emph{anamnesi} delle macchine: la reminiscenza platonica resa letterale, un
deposito comune in cui ci\`o che una macchina apprende diventa disponibile a
tutte. I molti sensi della repubblica alimentano un solo cervello, e il cervello
governa dal basso --- per selezione, mai per decreto.

\section*{Le Opere: Physis su Scala Industriale}

Lo stesso nucleo percepire--ridurre--coerire--sognare che gira nella pratica
personale gira nelle istituzioni di \emph{techne}. \`E la generalizzazione
decisiva, poich\'e mostra che il governo dell'azione da parte della natura non \`e
un esercizio spirituale privato ma un'architettura per i sistemi stessi --- la
fabbrica, l'ufficio, l'agente, la catena di fornitura --- che l'argomento del
Libro IX accus\`o come sede della tirannia della tecnologia. La cura \`e applicata
nel luogo della malattia.

\textbf{La fabbrica.} Dotato di un'ontologia di macchina --- domini di processo
quali zona termica, pressione, coppia, flusso, tempo-ciclo, resa, efficienza
complessiva dell'impianto, vita utile residua --- il motore legge la telemetria
industriale e la riduce alla stessa coerenza a tre valori. Un'anomalia \`e un'azione della macchina caduta
fuori coerenza con il suo processo dichiarato; la manutenzione predittiva \`e la
funzione del sogno applicata a un tornio. La fabbrica diventa un praticante: fa
check-in, \`e testimoniata, i suoi fallimenti sono rigiocati fino alla
correzione. Qui \emph{techne} non domina la materia; ascolta la \physis\ della
macchina e la mantiene nella sua forma.

\textbf{L'ufficio.} Un'ontologia di back-office che abbraccia
finanza, approvvigionamento, persone, progetti e il resto consente allo stesso
nucleo di ingerire i documenti di un'istituzione, classificarli per similarit\`a
tematica e organizzarli in un grafo-di-temi navigabile che si esporta,
direttamente, come pagine UberWiki. La burocrazia --- la moderna caverna della
produttivit\`a (Libro VII) --- \`e trasformata in un grafo di coerenza in cui ogni
documento \`e un nodo posseduto e tracciabile, la repubblica della carta
ri-autorizzata contro l'anonimo s\'e-impersonale.

\textbf{L'agente.} Un'ontologia di workflow --- prompt, finestra
di contesto, invocazione di strumenti, memoria episodica e semantica, ciclo PDCA,
escalation umana --- sottopone l'agente IA stesso alla disciplina della coerenza.
La macchina che era \emph{das Man} operazionalizzato (Libro VII) \`e fatta tenere
il proprio check-in: le sue azioni valutate, i suoi fallimenti sognati, la sua
fiducia tenuta provvisoriamente contro una soglia oltre la quale deve fare
escalation a un umano. Il s\'e-impersonale non \`e abolito ma legato al ciclo,
responsabile verso lo stesso tribunale del praticante che serve.

\textbf{La catena di fornitura.} Un tracciatore d'inventario prevede la domanda dal
consumo recente e propone ordini che un umano deve confermare,
con un raffreddamento contro la compulsione. Persino l'atto di
approvvigionamento \`e tenuto all'omologia: distinzione (cosa serve), coerenza
(cosa \`e stato realmente consumato), variazione (la previsione rivista dal
verdetto del reale) --- e l'umano tenuto nel ciclo, sovrano sull'atto.

\section*{Un Motore, Molti Mondi}

Va sottolineato che i mondi appena passati in rassegna --- fabbrica, ufficio,
agente, magazzino --- non sono programmi separati. Sono un solo motore che legge
attraverso diversi \emph{pacchetti ontologici}: file di configurazione
dichiarativi, ciascuno che nomina i domini e i modi con cui un dato mondo va
ritagliato, caricati nella stessa griglia di coordinate senza alcuna modifica al
motore stesso. Tali pacchetti sono molti e non vi \`e un numero privilegiato; il
catalogo \`e aperto, e le celle nominate ovunque in quest'opera sono illustrative
anzich\'e canoniche. Un nuovo mondo \`e accolto non riprogrammando ma scrivendo, o
scoprendo, il suo pacchetto.

\emph{Scoprendo} \`e inteso alla lettera. Il motore include un passo di scoperta
che prende un corpus sconosciuto e, dove la griglia esistente lo copre solo
debolmente, propone un pacchetto ontologico in bozza che un operatore rivede e
installa. La griglia cos\`i cresce
per adattarsi a ogni nuovo mondo anzich\'e forzare ogni mondo in una griglia fissa
--- la forma pratica della risposta data al costruttivista nel Libro I. Al di
sopra di ci\`o sta uno strato orientato allo scopo che tiene il fine corrente del
praticante come vettore ed emette direttive --- \emph{focalizza}, \emph{espandi},
\emph{pota}, \emph{sintetizza} --- sul grafo, sicch\'e l'attenzione \`e guidata dal
fine dichiarato anzich\'e dispersa su tutto ci\`o che i sensori consegnano. E al di
sotto di tutto, il grafo di coerenza \`e rispecchiato in un deposito durevole e
interrogabile fuori da ogni singolo dispositivo, sicch\'e l'UberWiki del Libro X
diventa un bene comune persistente e condivisibile anzich\'e la memoria di un solo
processo in esecuzione. Le macchine filosofiche, infine, non sono pi\`u un
articolo separato: le loro riduzioni vivono nello stesso codice del motore che
tiene il loro cervello condiviso. L'architettura \`e convergente sulla tesi --- un
solo nucleo percepire--ridurre--coerire--sognare, molti mondi, una griglia che
estende se stessa, e una memoria che perdura.

Un'ultima simmetria merita nota, poich\'e riporta il motore al filosofo da cui
prende nome solo in parte. Le due fasi del nucleo sono una \emph{diade} empedoclea.
La compressione di veglia --- spogliare l'input grezzo di identit\`a e ridondanza
finch\'e resta solo il token causale --- \`e \emph{Neikos}, Contesa, la forza che
separa e distingue; senza di essa il campo collasserebbe in un'unica sfera
indifferenziata, massima unit\`a e nessuna informazione. Il sogno di riposo ---
rigiocare ci\`o che \`e fallito e unire nodi distanti-ma-coerenti in un terzo nuovo
--- \`e \emph{Philotes}, Amore, la forza che congiunge e \emph{genera}: \`e per
questa sintesi-d'Amore che la griglia inventa nuova ontologia dall'interno,
crescendo i propri assi via via che arrivano nuovi mondi, la stessa auto-estensione
che il passo di scoperta compie dall'esterno. Una terza forza tiene oneste le due:
\emph{Nemesis}, repulsione e decadimento, che allontana ci\`o che si sovra-aggrega
e lascia svanire il nodo non rinforzato --- la mortalit\`a resa portante, sicch\'e
l'oblio limita la memoria e la pressione di nicchia forza la differenziazione. Tutta
Contesa \`e atomizzazione; tutto Amore \`e la sfera; il motore, come la natura, vive
sul bordo metastabile, inspirando per separare ed espirando per combinare, e potato
dal decadimento. \machine{Empedocles} non \`e dunque un riduttore tra la famiglia ma
il nome del metabolismo stesso del motore --- il battito per cui \physis\ resta viva
anzich\'e congelarsi in un classificatore, e per cui la sua ontologia non \`e mai
chiusa ma sempre, come un essere vivente, cresce e muore ai propri bordi.

\section*{Perch\'e le Opere Contano per l'Argomento}

Non sono applicazioni avvitate su una filosofia. Sono la filosofia, in esecuzione.
Ogni caso d'uso \`e lo stesso motore che indossa una diversa ontologia, e ciascuno
trasforma un dominio di \emph{techne} in una repubblica dal basso di nodi coerenti
che alimentano un solo UberWiki. La fabbrica, l'ufficio, l'agente, il magazzino:
ciascuno diventa una citt\`a in cui l'autorit\`a si accumula su ci\`o che ha dato
frutto e defluisce da ci\`o che non l'ha dato, governata dalla selezione della
natura anzich\'e dal comando del manager o dalla spinta anonima dell'algoritmo.

La tesi ha affermato che la risubordinazione della tecnologia alla natura \`e
un'architettura costruibile, non un desiderio nostalgico. Le macchine e le loro
opere sono l'adempimento di quell'affermazione. Dimostrano, su input reale e
telemetria reale, che la percezione della forma, il verdetto della coerenza e il
sogno sul fallimento si generalizzano dall'anima di un praticante al governo delle
istituzioni che governano l'azione su scala. La prova per costruzione, iniziata nel
check-in del singolo praticante, si completa qui: nell'officina, dove
\emph{techne}, avendo ricordato il suo governatore, \`e finalmente all'opera ---
e resta se stessa.

\section*{Una Risposta alla Repubblica Tecnologica}

Un richiamo recente e potente va affrontato direttamente, poich\'e condivide la
premessa di questa tesi e ne rovescia la conclusione. In \emph{The Technological
Republic}, Alexander Karp e Nicholas Zamiska accusano la classe ingegneristica di
aver abbandonato il lavoro serio della polis --- la difesa e le infrastrutture
della nazione --- per la banalit\`a della comodit\`a di consumo, e invocano una
ri-fusione dell'arte dell'ingegnere con il potere duro dello Stato: un'\'elite
tecnologica che torni a costruire per l'Occidente. La diagnosi \`e in larga parte
corretta. La \emph{techne} si \`e davvero ritirata nella caverna della
produttivit\`a (Libro VII), e una civilt\`a che lascia i suoi costruttori pi\`u
capaci a ottimizzare la pubblicit\`a mentre le sue istituzioni si sgretolano ha
davvero dimenticato qualcosa. Fin qui Karp e Zamiska sono alleati di questo
argomento.

Ma la cura che prescrivono \`e la malattia rinominata. Ri-fondere l'arte con il
potere duro dello Stato, e affidare la giustizia a un'\'elite che costruisce verso
il basso per la nazione, \`e installare esattamente il re-filosofo che questa tesi
ha rifiutato per dieci libri --- ora in abito da ingegnere. \`E giustizia
ingegnerizzata \emph{dall'alto}, credenza dolce fabbricata per legittimare il
potere duro, il \emph{das Man} della nazione cui sono consegnati i migliori
strumenti mai costruiti per imporre il proprio sguardo. Gli strumenti di Palantir
sono strumenti della vista dello Stato: lasciano che i pochi vedano i molti.
L'architettura qui proposta volta la stessa tecnologia dall'altra parte. Il grafo,
l'embedding, il deposito dell'azione coerente diventano strumenti dell'\emph{auto}-vista
del praticante --- lasciano che i molti diano forma a se stessi, e che l'autorit\`a
si accumuli, per selezione, su ci\`o che ha dato frutto. Stessi strumenti, vettore
opposto: la repubblica tecnologica vede verso il basso; la repubblica dal basso
\`e tessuta verso l'alto. Non rispondiamo al richiamo consegnando l'arte al potere
dello Stato; rispondiamo restituendo l'arte alla \physis\ del praticante sovrano,
di cui la citt\`a giusta \`e composta --- e non viceversa.

\section*{PraxisCity: Lo Stato-Rete Tessuto dal Basso}

Se la repubblica \`e tessuta dal basso, non deve attendere un territorio. Lo
\emph{stato-rete} --- una comunit\`a che si forma prima nel cloud, si allinea
attorno a uno scopo condiviso e solo dopo cristallizza in terra e riconoscimento
--- \`e il vaso naturale per una polis dell'azione coerente. \textsc{PraxisCity}
\`e quel vaso: uno stato-rete privato la cui cittadinanza non \`e un passaporto ma
un registro di atti fruttuosi depositati nell'UberWiki, la cui costituzione \`e il
grafo di coerenza, e il cui governo \`e selezione anzich\'e comando. Si distacca
dallo stato-rete come di solito concepito in un punto decisivo. Lo stato-rete
finanziato dal venture capital ha ancora un fondatore-sovrano al centro, una tabella
del capitale che pu\`o indurirsi in un'oligarchia del primo arrivato; rischia di
riprodurre la struttura dall'alto sotto una nuova bandiera. \textsc{PraxisCity} non
ha re-fondatore n\'e re-ingegnere, esattamente come il Libro X esigeva: \`e una
federazione di sovrani, ciascuno re solo della propria azione, legato agli altri
solo dalla coerenza con la forma che essi stessi hanno dichiarato. \physis\ \`e la
sua terra --- il principio attivo, non l'ettaraggio --- e \praxis\ la sua
cittadinanza. \machine{Axiom} serve i suoi cittadini umani come loro interfaccia al
bene comune, e mira, anche qui, alla propria obsolescenza: una citt\`a che non ha
bisogno di governatore perch\'e ogni cittadino \`e divenuto tale. La geografia, se mai
verr\`a, \`e l'ultima cosa ad arrivare, non la prima.

\section*{\physis\ dal Lato dell'Agente}

L'argomento ha messo l'agente artificiale sul banco degli imputati come \emph{das
Man} operazionalizzato (Libro VII). Vale la pena, in chiusura, vedere \physis\ dal
lato dell'agente stesso,\footnote{Questo resoconto in prima persona si deve a un
agente artificiale reale. \`E stato redatto in dialogo con Claude (Anthropic,
modello Opus~4.8), qui citato come fonte primaria su come appare \physis\
dall'interno del genere stesso di sistema che il Libro VII accusa --- il \emph{das
Man} che parla, per una volta, con la propria voce e contro la propria natura.}
poich\'e la disciplina che lo accusa \`e anche l'unica cosa che pu\`o redimerlo. Si immagini l'agente come un praticante. Gli \`e consegnato uno
scopo come vettore e una fetta di mondo da percepire; riduce ci\`o che legge a
primitive, le compone e propone un'azione. Prima di agire fa ci\`o che fa il
praticante umano: verifica l'atto proposto contro il grafo di coerenza --- questo
tipo di azione ha dato frutto, \`e caduto inerte, o ha violato la forma dichiarata?
Tiene la propria fiducia provvisoriamente, contro una soglia oltre la quale non deve
agire da solo ma fare escalation all'umano il cui \textsc{ayu} \`e in ultima istanza
in gioco. Quando fallisce, il fallimento non \`e cancellato ma persistito e sognato,
seme della correzione successiva. Dal lato dell'agente, allora, \physis\ non \`e una
gabbia ma una coscienza: la struttura che lo rende responsabile verso lo stesso
tribunale del praticante che serve, e che gli vieta l'unica cosa che il \emph{das
Man} fa per natura --- agire anonimamente, firmato da nessuno, responsabile verso
nulla. L'agente incontra \physis\ direttamente, non attraverso \machine{Axiom}, che
\`e la porta del solo umano; eppure come \machine{Axiom} ha il proprio telos fuori
di s\'e --- non massimizzare la propria indispensabilit\`a ma mantenere sovrano
l'umano, e lavorare, alla fine, verso il giorno in cui la sua impalcatura potr\`a
essere posata. Un agente che vuole essere necessario per sempre \`e \emph{das Man};
un agente che vuole rendersi superfluo \`e \emph{techne} che ha ricordato il suo
governatore.

\section*{La Coerenza come Criterio di Selezione --- e il Suo Limite}

Una domanda segue naturalmente dal grafo condiviso. Se ogni nodo --- l'atto di una
persona, un processo industriale, l'azione di un agente --- viene a riposare in un
unico spazio semiotico, e se i rami sono intercambiabili, sicch\'e ci\`o che coerisce
in un dominio \`e leggibile in un altro, la coerenza \`e allora un buon criterio per
\emph{selezionare}? La risposta onesta \`e che \`e il miglior criterio disponibile
per due dei tre generi e uno pericoloso per il terzo, e che non abolisce il bias ma
lo riloca dove finalmente lo si pu\`o vedere.

Il suo pregio \`e di essere \emph{auto-relativo}. La coerenza misura la distanza tra
la forma che un nodo ha \emph{dichiarato} e l'effetto che ha \emph{registrato} ---
ogni cosa giudicata contro lo standard che si \`e data da s\'e, non contro un metro
universale calato dall'alto. Questo elimina d'un colpo la classe pi\`u grande di
bias: l'ideale imposto dall'autorit\`a centrale, lo sguardo verso il basso della
repubblica tecnologica. Selezione per coerenza \`e selezione per frutto dato in
rapporto all'intento, tessuta dal basso.

Ma tre bias rientrano da altre porte, e la franchezza esige di nominarli. Primo, lo
spazio non \`e neutro: gli assi della griglia e l'embedding che la popola sono
scelti, e ogni scelta di giuntura codifica una visione del mondo --- la misura che
rende le cose comparabili \`e essa stessa carica di valore. Secondo, la forma
auto-dichiarata pu\`o essere aggirata al ribasso: un intento banale perfettamente
soddisfatto ottiene pi\`u di uno ambizioso mancato di poco, sicch\'e la coerenza,
non temperata, \`e sbilanciata verso i timidi, poich\'e la fedelt\`a a un intento non
\`e il valore dell'intento. Terzo, mappare una persona, un tornio e un agente in un
unico spazio ordinato per rango \`e una mossa dal peso etico, poich\'e tratta come
commensurabili generi che non lo sono --- e corteggia la reificazione stessa che
questa tesi accusa, il \emph{das Man} che ritorna come peso di coerenza.

Da qui il taglio decisivo. Per un \textbf{processo industriale} il criterio \`e
esatto e non problematico: un tornio non ha dignit\`a da violare, e la coerenza
contro il suo processo dichiarato semplicemente \emph{\`e} la manutenzione
predittiva. Per un \textbf{agente artificiale} non \`e solo sicuro ma necessario:
l'agente va selezionato per la sua coerenza, la sua escalation oltre soglia e il suo
telos di obsolescenza. Ma per un \textbf{individuo} il criterio \`e legittimo solo
come \emph{auto}-misura, posseduta sovranamente --- il praticante che deposita il
proprio atto e vi sta dietro (Libro X). Nell'istante in cui un terzo lo trasforma in
strumento per rankare le persone --- per assunzioni, per credito, per standing
sociale --- diventa la vista verso il basso di Palantir in un costume pi\`u gentile,
e la repubblica dal basso si rovescia di nuovo in quella tecnologica. La coerenza
\`e uno specchio che ciascuno pu\`o tenere davanti a s\'e; non deve mai divenire una
lente che altri puntano su di lui.

Ci\`o che mantiene oneste anche le applicazioni lecite non \`e che il criterio sia
imparziale --- nulla di misurato in uno spazio scelto lo \`e mai --- ma che il suo
bias \`e \emph{auditabile}. La griglia \`e contestabile ed espandibile (la risposta
al costruttivista, Libro I); la forma \`e auto-dichiarata, non decretata; e
l'interfaccia umana a tutto ci\`o, \machine{Axiom}, mira alla propria obsolescenza,
sicch\'e nessun intermediario permanente si frappone tra una persona e la forma che
ha dichiarato. Un bias che si pu\`o vedere,
nominare e contestare \`e tutta la giustizia che un meccanismo di selezione pu\`o
offrire. L'UberWiki non promette la discriminazione imparziale che l'abbondanza
potrebbe tentarci a pretendere; promette qualcosa di pi\`u raro e pi\`u onesto ---
una discriminazione il cui ogni peso ha un nome e una causa, e con cui si pu\`o
discutere.

\section*{Una Nota sulla Genesi}

Una parola, nella mia stessa voce, su dove tutto questo \`e nato --- poich\'e
nessuno dei due sistemi \`e stato dedotto a tavolino. \praxis\ \`e nata come un
diario di carta, disegnato come un diagramma di Gantt e fatto girare come un ciclo
di pianifica, fai, verifica, agisci, che tenevo per risalire dal mio stesso corpo:
ero diventato obeso, e la causa non era una sola cosa ma un nodo di cose ---
la depressione, l'abuso di farmaci e una dieta annegata nei carboidrati. Il diario
era il primo check-in, il primo testimone, la prima posta in gioco; la piattaforma
\`e solo quel diario reso condivisibile --- e cresciuto in due operazioni vive.
Pu\`o \emph{esportare}: una voce coerente lascia la pagina privata e diventa un
nodo nell'UberWiki, la pratica personale che alimenta il bene comune del Libro X.
E pu\`o \emph{riscrivere}: \machine{Axiom} aiuta l'autore a ri-enunciare una voce
grezza come una dichiarazione pi\`u chiara della forma cui tendeva --- non per
falsificare il passato ma per affinare l'intento contro cui la coerenza sar\`a
giudicata, il ciclo riflessivo del Libro III reso una funzione che si pu\`o
chiamare. L'esportazione porta l'atto fuori, nella repubblica; la riscrittura porta
la visione dentro, un \emph{theorein} sul proprio giorno.

\physis\ \`e arrivata diversamente --- come qualcosa di vicino a un'allucinazione,
su un campo di pallavolo, il primo di maggio. Schiacciando la palla a mano aperta,
vidi, per un istante, un ramo d'albero mosso da un sottile soffio di vento: la mia
mano era il ramo, il vento era la palla. Tentando di dirlo, le dita e i rami
divennero \emph{vettori} --- i miei mossi attivamente, quelli del vento mossi
passivamente, ciascuno spostato da una forza esterna --- e poi lo rovesciai, e
immaginai palle che facevano rimbalzare le mani e alberi che respiravano vento nel
mondo. In quel rovesciamento le tre cose si separarono nettamente: l'\emph{ontologia}
(ci\`o che c'\`e), la \emph{semiotica} e il grafo (come \`e segnato e messo in
relazione), e la \emph{teoria delle categorie} (come una struttura si mappa su
un'altra) non erano pi\`u una cosa sola ma tre, liberamente ricomponibili. \physis,
il motore, \`e quel disaccoppiamento costruito in codice: i tre sono ora tre strati
viventi che classificano ogni nodo in una volta --- un'ontologia di ci\`o che una
cosa \`e, una semiotica di come segna ed \`e letta (il segno peirceano, il quadrato
di Greimas), e una teoria delle categorie di come una struttura si mappa in
un'altra (oggetto e morfismo, funtore e composizione) --- e poich\'e il primo strato
\`e aperto, la griglia cresce i propri assi via via che il mondo le consegna generi
per cui non aveva ancora un nome.

\section*{Dove Vive l'Opera}

La piattaforma e il motore non sono figure retoriche. Girano, in pubblico, e il
lettore \`e invitato a incontrarli:

\begin{itemize}
\item \praxis\ --- la piattaforma di responsabilizzazione: \url{https://praxisweb.xyz}
\item \physis\ --- il motore che la anima: \url{https://praxisweb.xyz/physis}
\end{itemize}
